\section{Conclusion}\label{sec:conclusion}

I have studied how Peru's 2022 minimum wage increase, from 930 to 1,025 soles, affected
workers with different levels of education, using quarterly ENAHO data and a
difference-in-differences design that exploits the sharply different exposure of low- and
high-education workers to a single national reform. My central finding is that the reform
shifted the composition of low-skilled employment rather than the level of low-skilled
pay: formal employment fell by 4.6 percentage points, a fifth of its baseline, and
self-employment rose by 3.6 points, with flat formality pre-trends, larger shifts where
the floor cut deeper into the regional wage distribution, and confirmation from a
continuous-exposure design under randomization inference. No pay response accompanied the
shift: relative wages of low-education workers did not rise at the mean or anywhere in
the distribution, and the share of wage earners paid below the floor did not fall,
pointing to non-compliance as the reason the floor could not lift pay. The overall
employment effect cannot be pinned down and is confounded by the uneven post-pandemic
recovery. Around the identically sized reform of January 2025, in a calmer labor market,
neither wages nor composition responded, which suggests that the reallocation margin
operates when the labor market holds many marginal formal matches, as it did during the
post-pandemic re-sorting of 2022, and lies dormant when it does not.

The most economical reading of these results is that, in a labor market where informality
is pervasive and enforcement is weak, a moderate and inflation-eroded increase in the
statutory minimum does not lift the pay of the low-skilled; what it moves, when it moves
anything, is the allocation of workers between formal and informal arrangements. This
qualifies the lighthouse view of the Peruvian minimum wage that emerges from studies of
earlier reforms and points to the informal margin as the central channel of adjustment.

Two directions for future research follow. The first is to exploit panel data, either
the panel component of ENAHO or the higher-frequency EPEN, to follow individuals across
reforms and distinguish worker-level transitions out of formality from changes in the
composition of hiring, which my repeated cross-sections cannot separate. The second is to
study the interaction between minimum wage policy and enforcement directly, since my
results suggest that the returns to raising the floor depend on whether the floor can be
made to bind. As additional post-reform quarters accumulate, it will also become possible
to trace the longer-run adjustment that a short window such as mine cannot capture.
